\section{Introduction}
\label{sec:intro}

% Local fallbacks for macros used in the overview figures (kept here to avoid touching the preamble).

Anonymous credentials (AC) are cryptographic protocols that allow parties to convey certified information about themselves in a privacy-preserving way. There are usually three roles in ACs, namely \emph{issuers}, \emph{holders} and \emph{verifiers}.  More concretely, an AC scheme allows a holder to prove to a verifier that it possesses certain attributes, which have been authenticated by the issuer. Here, the authenticated attributes are usually called a \emph{credential}, which is kept private by the user. Using a so-called \emph{showing} process, the user can generate a presentation which convinces the verifier that attributes stored in its credential (or credentials) satisfy predicates without revealing any additional information. While there are numerous security properties that ACs can achieve, the most fundamental ones are unforgeability and unlinkability. Here, the former ensures that an issuer indeed generated a credential that a user presents, and that colluding issuers and verifiers may not link a presentation to a specific issuance session.  Here, the latter should also ensure that multiple presentations cannot be traced to the same credential.

Certain AC use-cases like rate limiting~\cite{yun-privacypass-crypto-arc-00,schlesinger-cfrg-act-00} require that the credential is only valid for a \emph{bounded} number of presentations in a given context (e.g. per epoch or per service) -- without revealing which presentations come from the same holder up to the enforced limit. 
A key advantage of existing, pre-quantum ACs is their performance, e.g. with constant credential and presentation size, regardless of the enforced rate limit. However, the situation looks very different in the post-quantum setting. There, practically efficient AC schemes~\cite{CCS:BLNS23,CCS:AGJLRS24,EPRINT:PWFW23}, are yet mostly absent. This is highly unfortunate, given the uptake of (privacy-preserving) digital identification schemes in e.g. the EU Digital Identity Wallet, and simultaneous requirement for deployed schemes to eventually be quantum-secure: given current migration timelines, post-quantum ACs may become mandatory in the next 5 years or fewer. Recently, Google~\cite{google-pq} moved the timeline for post-quantum to 2029. 
While everlasting anonymous rate-limited tokens~\cite{AC:CDLT25} provide post-quantum privacy, they only guarantee rate limiting until quantum computers arrive. 



\paragraph*{Building ACs.}
To understand how post-quantum ACs (with rate-limiting) are usually  constructed, it is instructive to understand how  classically-secure constructions work. 
Anonymous credentials are often built on top of a Blind Signature design. Here, one of the most well-known approach (which is also the most promising in the post-quantum setting) is due to Fischlin~\cite{C:Fischlin06}:
\begin{enumerate}
	\item The issuer initially creates a key pair $(\sk,\vk)$ using an EUF CMA-secure signature scheme. We assume that user and verifier have $\vk$, while the issuer keeps $\sk$.
	\item The holder first commits to its message $\mu$ as \(t\), using randomness $r$. It then sends $t$ to the issuer.
	\item Upon receiving $t$, the issuer signs it and thus  creates a signature $\sig$. It then sends $\sig$ to the user, who checks its verifiability with $\vk$.
	\item To create the blind signature, the user creates a Non-Interactive Zero-Knowledge Proof of Knowledge (NIZK) showing that it knows $\sig,r$ such that $$\Verify_{\vk}(t,\sig)=1 \land t= \Commit(\mu;r)$$
\end{enumerate}
Combining existing post-quantum primitives generically, even with zero-knowledge friendly hash functions~\cite{EPRINT:PWFW23}, leads to non-competitive performance. Thus, existing works carefully choose commitment schemes, signatures and NIZKs such that the NIZK for the given relation has good proof size and prover/verifier time. However, the state-of-the-art is generally not considered efficient enough for practical deployments:
Constructing blind signatures from alternative post-quantum hash-and-sign
cryptography~\cite{EPRINT:BFMRSV25a} and anonymous credentials from
lattices~\cite{EPRINT:LCLLW23}, where recent work brings the credential size
down to just over 70 kB using optimized samplers~\cite{C:JeuRouSan23,AC:JeuSan25}. 

Our construction builds on proof-friendly lattice signatures from vanishing SIS (vSIS) with rational hashing
\cite{EPRINT:DKLW25,PKC:DKLW25} and on fast trapdoor generation and Gaussian preimage sampling for NTRU-like
lattices \cite{EPRINT:ENSTW23}. For succinct zero knowledge over small fields, we use the SmallWood framework
for hash-based polynomial commitments and compiled PIOPs \cite{EPRINT:FenRiv25b}. 
While SmallWood has been used for ZK-friendly post-quantum signatures in CAPSS~\cite{EPRINT:FenRiv25a}, we use a zk-friendly PRF with a significantly smaller modulus for tag derivation, which reduces the size of our communication significantly, and discuss Poseidon/Poseidon2-style instantiations \cite{USENIX:GKRRS21,AFRICACRYPT:GraKhoSch23}.



\subsection{Contributions}
In this work, we narrow the efficiency gap for post-quantum blind signatures and ACs substantially. For Blind Signatures, we achieve this by optimizing existing design ideas carefully and employ novel primitives. We then extend our construction to yield the first post-quantum rate-limiting anonymous credentials. Crucially, the size of our credential as well as the presentation is essentially independent of the rate limit.

We construct a Blind Signature based on a signature whose security relies on lattice assumptions together with a  hash-based NIZK. We then extend this to an \emph{Anonymous Rate-Limited Credential}~(ARC)~\cite{yun-privacypass-crypto-arc-00} by including a Pseudorandom Function (PRF) into our design. Using a key committed in the message that is signed by the issuer, the user is able to evaluate a PRF on the key on different nonces. This PRF output, which we call a \emph{tag}, can be made public. Since the signature and thus the key remain hidden due to the use of a NIZK,  this intuitively makes the scheme unlinkable. By including a proof of correct PRF evaluation into the NIZK, the user also cannot generate a differing tag for the same nonce. Therefore, any two presentations generated for the same credential are linkable as they must have the same tag, which makes the scheme rate-limiting.

To design the Blind Signature, we follow the recent proposal of \cite{C:BLNS23,CCS:LyuSeiSte24,PKC:DKLW25} which also presents blind signatures and anonymous credentials based on lattice assumptions. We first let the holder commit to its message $\mu$ using an Ajtai commitment. Then, in an interactive process, issuer and signer will sample randomness which together with the committed message becomes  a target \(\vect\)  that will be signed by the issuer. Here, the user proves in zero knowledge that \(\vect\) is consistent with the committed message and the issuer-supplied randomness. The signature scheme is formed by a matrix $\mA$ together with a trapdoor $\trapdoor$ that allows sampling of short preimages $\bmod q$ using a GPV-style~\cite{STOC:GenPeiVai08} preimage sampler.
	The issuer then uses $\trapdoor$ to sample a short preimage \(\vecu\) such that \(\mA \vecu=\vect \bmod q\), which serves as the signature. To prove knowledge of a valid $\vecu$ we then use the SmallWood NIZK scheme~\cite{EPRINT:FenRiv25b}. While this approach makes the round complexity of issuance non-optimal (requiring 3 rounds instead of 1 round), it allows tighter parameters for the preimage sampler (as $\vect$ is uniform) while making the statement proven in the NIZK structurally simple. To turn this construction into an ARC, we let the user additionally commit to a PRF key $\veck$ for the Poseidon2 hash function~\cite{AFRICACRYPT:GraKhoSch23}. It turns out that the parameters such as the modulus $q$ of Poseidon2 can be chosen to be compatible with those of the lattice commitment, signature scheme and NIZK, making the proof of PRF evaluation highly efficient. The holder can then attach an evaluation of the PRF on a nonce where the committed key is used, together with a proof of correct PRF evaluation, as part of the presentation. If the holder reuses the nonce, this will be noticed by a verifier keeping track of previously issued tags.



\subsection{Technical overview}

This work constructs ARC around a single protocol blueprint with message  \(\vecm=\vecmu\Vert \veck\), where \(\vecmu\) is a vector of
attributes and \(\veck\) is a \emph{secret} PRF key.

\paragraph*{Issuance}
Let $\mACom$ be a uniformly random and compressing matrix to be used for the Ajtai commitment. In the issuance protocol, the holder additionally samples  short $\vecr_0^\Holder,\vecr_1^\Holder,\vecr$  to compute the commitment $\vecc=\mACom [\vecmu \Vert \vecr]$ that is sent to the issuer. It, in response, samples short $\vecr_0^\Issuer,\vecr_1^\Issuer$ that are returned to the holder. The target $\vect$ is then computed as
	$$ \vect :=hash(\vecmu,\veck,\centerfunc{\vecr_0^\Holder+\vecr_0^\Issuer},\centerfunc{\vecr_1^\Holder+\vecr_1^\Issuer})$$
	where $\centerfunc$ derives a short value from the sum and $hash$ is a
  so-called rational hash function as introduced in~\cite{C:BLNS23}. The holder
  then sends $\vect$ to the issuer, together with a NIZK that it is well-formed according to $\vecc,\vecr_0^\Issuer,\vecr_1^\Issuer$ and that all involved values are short.

 The signature scheme is based on a hash-and sign signature. It precomputes preimages of a matrix $\mA$ using a trapdoor Gaussian preimage sampler over an
	NTRU-like lattice (Antrag/hybrid sampling) to instantiate the issuer’s trapdoor signing step. The short preimage $\vecu$ can then be returned to the holder, who verifies that $\vect = \mA \vecu$.
	
	We illustrate the issuance transcript flow and pre-sign proofs in~\Cref{fig:issuance-flow}.
	

\begin{figure}[H]
	\centering
    \resizebox{\columnwidth}{!}{%
	\begin{tikzpicture}[
		actor/.style={font=\bfseries, align=center},
		lifeline/.style={draw, thick},
		msg/.style={-Latex, thick},
		note/.style={draw, rounded corners, fill=gray!10, align=center, inner sep=4pt},
		x=1.15cm, y=0.9cm
		]
		% --- Actors ---
		\node[actor] (I) at (0,10)  {\Large Issuer\\[-1pt]$\Issuer$};
		\node[actor] (H) at (7,10)  {\Large Holder\\[-1pt]$\Holder$};
		
		% --- Lifelines ---
		\draw[lifeline] (I) -- ++(0,-9.5);
		\draw[lifeline] (H) -- ++(0,-9.5);
		
		% --- Notes and messages ---
		\node[note] at ($(H)+(0,-1.55)$) {Sample $\vecm=\vecmu\Vert \veck$\\ (with $\veck$ as PRF key)\\ Sample $\vecr_0^\Holder,\vecr_1^\Holder,\bar \vecr$};
		
		% Commitment (Holder -> Issuer)
		\draw[msg] ($(H)+(0,-3.2)$) -- node[above,sloped]
		{$\vecc=\mACom(\vecmu; \vecr_0^\Holder,\vecr_1^\Holder,\bar \vecr)$} ($(I)+(0,-3.2)$);
		
		\node[note] at ($(I)+(0,-4.0)$) {Sample issuer challenges\\ $\vecr_0^\Issuer,\vecr_1^\Issuer \getsunif [-\BoundMsg,\BoundMsg]$};
		
		% Challenges (Issuer -> Holder)
		\draw[msg] ($(I)+(0,-4.8)$) -- node[above,sloped] {$\vecr_0^\Issuer,\vecr_1^\Issuer$} ($(H)+(0,-4.8)$);
		
		\node[note] at ($(H)+(0,-5.8)$) {Center-combine\\ $\vecr_b \leftarrow \centerfunc(\vecr_b^\Holder+\vecr_b^\Issuer)$\\ for $b\in\{0,1\}$};
		
		% Pre-sign target + proof (Holder -> Issuer)
		\draw[msg] ($(H)+(0,-6.8)$) -- node[above,sloped]
		{$\vect=h_{(\vecmu\Vert \veck),(\vecr_0,\vecr_1)}(B),\;\pi_{\mathsf{pre}}$} ($(I)+(0,-6.8)$);
		
		\node[note] at ($(I)+(0,-7.8)$) {Trapdoor sample\\ $\vecu\gets \SampPre(\trapdoor,\vect)$\\ such that $\mA \vecu = \vect$};
		
		% Signature (Issuer -> Holder)
		\draw[msg] ($(I)+(0,-8.8)$) -- node[above,sloped] {Signature $\vecu$} ($(H)+(0,-8.8)$);
		
		\node[note] at ($(H)+(0,-9.6)$) {Store credential\\ $(\vecmu,\veck,\vecr_0,\vecr_1,\vecu)$};
		
	\end{tikzpicture}
    }
	\caption{Issuance flow (blind signature). The pre-sign proof $\pi_{\mathsf{pre}}$ binds $\vect$ to the commitment and the centered randomness, while hiding $\mu$.}
	\label{fig:issuance-flow}
\end{figure}

	
	
\paragraph{Blind Signatures and ARC Presentations.}
Let $\vecr_0:=\centerfunc(\vecr_0^\Holder+\vecr_0^\Issuer),\vecr_1:=\centerfunc(\vecr_1^\Holder+\vecr_1^\Issuer)$. A presentation is a NIZK showing knowledge of $\vecm$, $\vecr_0,\vecr_1$, $\vect$,  such that  $\vect=hash(\vecm,\vecr_0,\vecr_1)$, as well as $\vecu$ such that $\vect =\mA \vecu$ where $\vecm,\vecr_0,\vecr_1,\vecu$ are short. Towards constructing ARC, we additionally compute $\prftag=\PRF(\veck;\nonce)$, append it to the NIZK and additionally prove in the NIZK that $\nonce$ is in the rate limiting interval and that $\prftag$ is well-formed.
	

The verifier checks the NIZK proof and enforces the rate policy by maintaining server-side state keyed by
\((\nonce,\prftag)\) (or an application-specific encoding thereof) to avoid double-spending.

\begin{figure}[H]
	\centering
    \resizebox{\columnwidth}{!}{%
	\begin{tikzpicture}[
		actor/.style={font=\bfseries, align=center},
		lifeline/.style={draw, thick},
		msg/.style={-Latex, thick},
		note/.style={draw, rounded corners, fill=gray!10, align=center, inner sep=4pt},
		x=1.15cm, y=0.9cm
		]
		% --- Actors ---
		\node[actor] (H) at (0,0)  {\Large Holder\\[-1pt]$\Holder$};
		\node[actor] (V) at (9,0)  {\Large Verifier\\[-1pt]$\Verifier$};
		
		% --- Lifelines ---
		\draw[lifeline] (H) -- ++(0,-7.9);
		\draw[lifeline] (V) -- ++(0,-7.9);
		
		% --- Notes and messages ---
		\node[note] at ($(H)+(0,-2.5)$) {Choose unused $\nonce$\\ Compute $\prftag\leftarrow \PRF(\veck,\nonce)$ \\ Compute $\pi_{\mathsf{post}}$ proves:\\ $\exists(\vecu,\vecmu,\veck,\vecr_0,\vecr_1)$ s.t.\\ $\mA \vecu=h_{(\vecmu\Vert \veck),(\vecr_0,\vecr_1)}(B)$\\ and $\prftag=\PRF(\veck,\nonce)$
		};
		
		% Presentation (Holder -> Verifier)
		\draw[msg] ($(H)+(0,-4.3)$) -- node[above,sloped]
		{$(\prftag,\nonce,\pi_{\mathsf{post}})$} ($(V)+(0,-4.3)$);
		
		\node[note] at ($(V)+(0,-5.9)$) {Stateful rate check:\\ reject if $(\nonce,\prftag)$ already seen \\ else store $(\nonce,\prftag)$\\ reject if $\Verify(\pi_{\mathsf{post}})$ fails \\ else accept};
		
		
		
		
		% Result (Verifier -> Holder)
		\draw[msg] ($(V)+(0,-7.5)$) -- node[above,sloped] {accept/reject} ($(H)+(0,-7.5)$);
		
	\end{tikzpicture}
        }
	\caption{Showing (ARC presentation). Using an unused $\nonce\leq
		\mathsf{presentation_limit}$, the holder computes a deterministic $\prftag$
		from $(\veck,\nonce)$ (rate limiting) and pseudorandom across unused nonces (unlinkability).}
	\label{fig:showing-flow}
\end{figure}

Figure~\ref{fig:showing-flow} summarizes how tags enable stateful rate limiting without revealing the credential.



\paragraph{SmallWood Proofs.}

SmallWood~\cite{EPRINT:FenRiv25b} is a hash-based polynomial commitment scheme
which produces efficient witnesses for polynomials of small degree $\leq
2^{16}$, which is particularly useful for
lattice-based problems. The construction avoids the overhead of
asymptotically-designed systems like STARKS while achieving better concrete
efficiency than MPC-in-the-head variants for small witnesses. SmallWood exclusively relies on
hash-based primitives, making it plausibly post-quantum secure. 
We use SmallWood as an 
NIZK for polynomial constraints over a small field \(\Fq\), and we treat it as a black box that provides:
(i) commitment to witness rows packed on an evaluation domain, (ii) succinct openings at verifier-chosen points, and
(iii) soundness and zero knowledge after Fiat--Shamir \cite{EPRINT:FenRiv25b}.

To obtain compact proofs, both issuance and showing are compiled into SmallWood statements over committed witness rows, with explicit bridge families for the linear commitment, centering, rational-hash relation, signature relation, and PRF evaluation. The exact compiled relations are specified in Section~\ref{sec:smallwood-model}, while proof-system internals are deferred to Appendix~\ref{app:smallwood}.

In the showing proof, the rational-hash/signature relation is realized through a full transform-basis replay image whose vanishing implies the corresponding coefficient-domain cleared relation by the theorem stated in Section~\ref{sec:smallwood-model}.
